\begin{frame}{How do we test sublinearity numerically?}
  \framesubtitle{Q2}
  \vspace{0.10cm}
  \begin{columns}[T,onlytextwidth]
    \begin{column}{0.31\textwidth}
      \textbf{1. Normalised regret}
      \[
        \widehat\rho_T=\frac{\widehat R_T}{T}
      \]
      \smallnote{Sublinear behaviour requires \(\widehat\rho_T\) to trend toward zero as the horizon grows.}
    \end{column}
    \begin{column}{0.31\textwidth}
      \textbf{2. Incremental slope}
      \[
        \widehat s_{i}=\frac{\widehat R_{T_{i+1}}-\widehat R_{T_i}}
        {T_{i+1}-T_i}
      \]
      \smallnote{A falling late-horizon slope indicates decreasing marginal regret.}
    \end{column}
    \begin{column}{0.31\textwidth}
      \textbf{3. Effective exponent}
      \[
        \widehat\alpha_i=
        \frac{\log \widehat R_{T_{i+1}}-\log \widehat R_{T_i}}
        {\log T_{i+1}-\log T_i}
      \]
      \smallnote{Values below one are consistent with \(\widehat R_T\propto T^{\alpha}\) over that interval.}
    \end{column}
  \end{columns}
  \vspace{0.42cm}
  \flatbox{
    \textbf{Protocol.} Evaluate a common set of horizons
    \(\{10^3,3\!\times\!10^3,10^4,3\!\times\!10^4,10^5\}\), average over 500 reruns, and examine all three diagnostics together.
  }
  \vspace{0.08cm}
  \micro{These are finite-horizon diagnostics. They provide evidence consistent with or against sublinearity; they do not prove an asymptotic theorem.}
\end{frame}
